\scanpage{45}\mspage{32}

Alors la flèche $u_{Z_1}$ se déduit de $u_{Z_0}$ par cochangement de base, de façon plus précise on a \deleted{\unsure{On}}un carré cocartésien, où

\diagram{\begin{tikzcd}[ampersand replacement=\&,column sep=small,row sep=small]
T={}\&X_0\arrow[r]\arrow[d]\&X_1\arrow[d]\\
\&Y_0\arrow[r,dash]\&Y_1
\end{tikzcd}}

et

\diagram{\begin{tikzcd}[ampersand replacement=\&,column sep=small,row sep=small]
T'={}\&X'_0\arrow[r,dash]\arrow[d]\&X'_1\arrow[d]\\
\&Y'_0\arrow[r]\&Y'_1
\end{tikzcd}}

désignent les faces \unsure{postérieure|postérieures} et \unsure{antérieure|antérieures} du cube (2.6.1) :

\diagram{\begin{tikzcd}[ampersand replacement=\&]
\&\widetilde Z_0\arrow[r,"u_{Z_0}"]\arrow[d]\arrow[dr,phantom,"\text{cocart}"]\&Y'_0\arrow[d]\\
(\unsure{*|\alpha})\&\widetilde Z_1\arrow[r,"u_{Z_1}"]\&Y'_1,
\end{tikzcd}}

\deleted{déduit du diagramme \unsure{multicocartésien}} qu'on voit sur la base du diagramme

\diagram{\begin{tikzcd}[ampersand replacement=\&,row sep=small,column sep=small]
\&X_0\arrow[rr]\arrow[dr,dash]\arrow[dd,dash,"f_0"']\&\&X_1\arrow[dr,dash]\arrow[dd,"f_1",near end]\&\&\\
\&\&X'_0\arrow[rr,dash,crossing over]\arrow[dd,dash]\arrow[dddr,"f^{\unsure{\prime}}_0",near start]\&\&X'_1\arrow[dd]\arrow[dddr,"f'_1"]\&\\
(\unsure{**|\alpha\alpha})\&Y_0\arrow[rr]\arrow[dr]\&\&Y_1\arrow[dr]\&\&\\
\&\&\widetilde Z_0\arrow[rr,dash]\arrow[dr,shift left=.5]\arrow[dr,shift right=.5,"u_{Z_0}"']\arrow[ur,phantom,"\text{cocart}",sloped]\&\&\widetilde Z_1\arrow[dr,"u_{Z_1}"']\&\\
\&\&\&Y'_0\arrow[rr]\arrow[ur,phantom,"\text{cocart}",sloped]\&\&Y'_1,
\end{tikzcd}}

dans lequel le cube est un cube totalement cocartésien, déduit du carré cocartésien $X_\bullet$ en le ``\unsure{copoussant}'' par la flèche $f_0:X_0\longrightarrow Y_0$. La base du cube, et le carré

% archive sheet number: 44

